\documentclass[11pt]{amsart}
\usepackage[T1]{fontenc}
\usepackage{lmodern}
\usepackage{amsmath,amssymb,amsthm,mathtools}
\usepackage{xcolor}
\usepackage[colorlinks=true,linkcolor=blue!55!black,citecolor=blue!55!black,urlcolor=blue!55!black]{hyperref}
\newtheorem{theorem}{Theorem}[section]
\newtheorem{lemma}[theorem]{Lemma}
\newtheorem{proposition}[theorem]{Proposition}
\theoremstyle{definition}
\newtheorem{definition}[theorem]{Definition}
\title{The Strict-Pair Product Extremal Value at 62 Is 31}
\author{Zijian Zeng}
\address{Institute of Computer Science and Digital Innovation, UCSI University, Kuala Lumpur, 56000, MALAYSIA}
\email{zijianzeng@foxmail.com; 1002266693@ucsiuniversity.edu.my}
\date{August 31, 2026}
\hypersetup{pdftitle={The Strict-Pair Product Extremal Value at 62 Is 31},pdfauthor={Zijian Zeng}}
\begin{document}
\begin{abstract}
Let $F(N)$ be the largest cardinality of a subset $A\subseteq[N]$ for which
the products $ab$, over two-element subsets $\{a,b\}\subseteq A$, are all
distinct.  We certify the finite equality $F(62)=31$.  The lower bound is an
explicit set whose 465 products are distinct.  For the upper bound,
equal-product collisions define a four-uniform hypergraph with 1203 edges.
An exact exhaustive vertex-cover enumeration rejects every cover of size at
most 30.  This finite result does not settle the asymptotic questions
associated with Erd\H{o}s problem 425.
\end{abstract}
\maketitle

\section{Introduction}
Erd\H{o}s studied sets of positive integers having distinct products of two
different members in 1938~\cite{Erdos1938}.  The classical theory determines
the order of the secondary term beyond $\pi(N)$ but not its limiting
coefficient; see Erd\H{o}s's graph-theoretic account~\cite{Erdos1969} and the
modern summary in Liu and Pach~\cite{LiuPach2019}.  Here we address only one
finite instance of the strict-pair problem.

The qualification ``strict-pair'' is essential.  We compare products indexed
by two-element subsets and therefore do not compare squares $a^2$ with other
products.  Some uses of ``multiplicative Sidon'' impose the stronger condition
that includes repeated factors.  No such strengthening is used below.

\section{Collision hypergraph}
Write $[N]=\{1,\ldots,N\}$.
\begin{definition}
A set $A\subseteq[N]$ is \emph{strict-pair product injective} if, for all
$a,b,c,d\in A$,
\[
 (a<b\ \wedge\ c<d\ \wedge\ ab=cd)
 \quad\Longrightarrow\quad (a=c\ \wedge\ b=d).
\]
Let $F(N)$ be the maximum cardinality of such a set.
\end{definition}

\begin{lemma}\label{lem:four}
If two distinct strict pairs of positive integers have the same product, then
their union has four elements.
\end{lemma}
\begin{proof}
If the pairs shared an element, cancellation of that positive element would
make their other elements equal.  The pairs would then be identical.
\end{proof}

Let $H_N$ be the four-uniform hypergraph on $[N]$ whose edges are the
four-sets that admit a partition into two pairs with equal products.
\begin{proposition}\label{prop:reduction}
For every $N\geq1$,
\[
                         F(N)=\alpha(H_N)=N-\tau(H_N),
\]
where $\alpha$ and $\tau$ denote the independence and vertex-cover numbers.
\end{proposition}
\begin{proof}
By Lemma~\ref{lem:four}, a set violates strict-pair product injectivity if and
only if it contains an edge of $H_N$.  Thus admissible sets are precisely the
independent sets.  Taking complements exchanges independent sets and vertex
covers in every finite hypergraph.
\end{proof}

\section{The exact finite value}
\begin{theorem}\label{thm:main}
Under the strict-pair convention, $F(62)=31$.
\end{theorem}
\begin{proof}
For the lower bound, take
\begin{align*}
A=\{&2,5,9,11,19,21,31,32,34,35,37,39,41,43,44,46,47,49,50,51,\\
    &52,53,54,55,56,57,58,59,60,61,62\}.
\end{align*}
Direct integer multiplication gives 465 distinct products from its 465
two-element subsets.  Hence $F(62)\geq31$.

For the reverse inequality, the exact enumeration described in
Section~\ref{sec:certificate} constructs $H_{62}$ in two independent ways;
both give the same 1203 edges.  It exhaustively rejects a vertex cover of size
at most 30.  Thus $\tau(H_{62})\geq31$, and
Proposition~\ref{prop:reduction} gives $F(62)\leq31$.
\end{proof}

\section{Exact enumeration certificate}\label{sec:certificate}
The instance uses no floating-point arithmetic.  One edge generator groups
all $\binom{62}{2}$ pairs by their integer product and takes the four-element
supports of distinct pairs in a common group.  A second generator examines
every four-subset and its three pairings.  The sorted edge lists agree exactly.

The cover procedure maintains the uncovered edges, the vertices still
available, and the remaining budget.  Singleton candidate sets are forced.
Two lower bounds are used: the number of uncovered edges divided by the
largest current vertex degree, rounded upward, and the size of a greedily
constructed family of edges with pairwise disjoint available candidate sets.
Both are lower bounds for every completion.

For branching, choose an uncovered edge, and denote its available candidates
by $v_1,\ldots,v_t$.  Branch $i$ selects $v_i$ and excludes
$v_1,\ldots,v_{i-1}$.  Every cover has a unique first selected candidate in
this order, so the branches are disjoint and exhaustive.  Only complete
states already proved infeasible are memoized.  The budget-30 run visits
4,684,599 recursive states and rejects the budget; an independent replay
gives the same state count.

Literal subset enumeration agrees with the cover procedure for every
$N\leq12$, and the two edge generators agree at $N=62$.  The archived verifier
reconstructs the hypergraph from the definition instead of trusting the
serialized instance.

\section{Scope and data availability}
The result is an isolated certified finite equality.  It neither determines
an asymptotic coefficient nor answers the higher-product question in
Erd\H{o}s problem 425.  The accompanying source code, serialized instance,
exact witness, run logs, small-instance cross-check, and release manifest
permit reconstruction of the complete enumeration using exact integer
arithmetic.

\bibliographystyle{amsplain}
\bibliography{references}
\end{document}
